\documentclass[11pt, aspectratio=169]{beamer}

% Gọi file cấu hình chung
\usepackage{./common_style}
\usepackage{tikz} % Thêm thư viện để vẽ hình
\usetikzlibrary{arrows.meta, calc} % Thư viện để tùy chỉnh mũi tên và tính toán tọa độ

\title[LOGARIT (LOGARITHM)]{HÀM SỐ MŨ VÀ HÀM SỐ LOGARIT}

\author[Trần Đức Mạnh]{
    Trần Đức Mạnh \\[0.2cm] 
}

\titlegraphic{}

\begin{document}

% ==================== TRANG BÌA ====================
\begin{frame}
    \titlepage
\end{frame}

\section{Khái niệm Logarit}

\subsection{Định nghĩa Logarit}

\begin{frame}[fragile]
    \frametitle{\thesection. \insertsection}
    \framesubtitle{\thesection.\thesubsection. Định nghĩa (Definition)}

    \begin{columns}[T, onlytextwidth]
        \begin{column}{0.48\textwidth}
            \begin{block}{Định nghĩa}
                \footnotesize
                Cho hai số thực dương $a$ và $b$ với $a \neq 1$.
                
                Số thực $\alpha$ thỏa mãn đẳng thức $a^\alpha = b$ được gọi là \textbf{logarit cơ số $a$ của $b$}, kí hiệu là $\log_a b$.
                
                \[
                    \boxed{\alpha = \log_a b \iff a^\alpha = b}
                \]
            \end{block}
            
            \vspace{0.1cm}
            
            \begin{block}{Tính chất cơ bản}
                \footnotesize
                Với $0 < a \neq 1$ và $b > 0$, ta có:
                \begin{itemize}
                    \item $\log_a 1 = 0 \quad$ (vì $a^0 = 1$)
                    \item $\log_a a = 1 \quad$ (vì $a^1 = a$)
                    \item $a^{\log_a b} = b$
                    \item $\log_a (a^\alpha) = \alpha$
                \end{itemize}
            \end{block}
            
            \vspace{0.1cm}
            \begin{alertblock}{Điều kiện tồn tại}
                \footnotesize
                Biểu thức $\log_a b$ chỉ có nghĩa khi: \\
                Cơ số $\mathbf{a > 0, a \neq 1}$ và Biểu thức dưới logarit $\mathbf{b > 0}$.
            \end{alertblock}
        \end{column}
        
        \begin{column}{0.48\textwidth}
            \centering
            \vspace{0.1cm}
            \begin{tikzpicture}[>=Stealth, thick, scale=0.75, every node/.style={scale=0.85}]
                % Trục tọa độ
                \draw[->, gray!60, thin] (-1,0) -- (5.5,0) node[right, text=black] {$x$};
                \draw[->, gray!60, thin] (0,-1) -- (0,5.5) node[above, text=black] {$y$};
                
                % Đường phân giác y = x
                \draw[dashed, gray] (-1,-1) -- (5,5) node[above right, text=black] {$y=x$};
                
                % Đồ thị hàm số mũ y = 2^x (Màu xanh)
                \draw[blue, ultra thick, domain=-1:2.3, samples=100] plot (\x, {2^\x}) node[left] {$y = a^x$};
                
                % Đồ thị hàm số logarit y = log_2(x) (Màu đỏ)
                \draw[red, ultra thick, domain=0.4:5.2, samples=100] plot (\x, {ln(\x)/ln(2)}) node[below right] {$y = \log_a x$};
                
                % Thể hiện sự đối xứng (tương quan a^\alpha = b và \log_a b = \alpha)
                \coordinate (A) at (2,4);   % Điểm trên y = 2^x
                \coordinate (B) at (4,2);   % Điểm trên y = log_2(x)
                
                % Đường gióng cho điểm A
                \draw[dashed, orange, thin] (2,0) node[below, text=black]{$\alpha$} -- (A);
                \draw[dashed, orange, thin] (0,4) node[left, text=black]{$b$} -- (A);
                
                % Đường gióng cho điểm B
                \draw[dashed, green!60!black, thin] (4,0) node[below, text=black]{$b$} -- (B);
                \draw[dashed, green!60!black, thin] (0,2) node[left, text=black]{$\alpha$} -- (B);
                
                % Kẻ đường nối A và B vuông góc với y=x
                \draw[gray, dotted, thick] (A) -- (B);
                
                % Điểm M và M'
                \filldraw[black] (A) circle (1.5pt) node[above left] {$M(\alpha; b)$};
                \filldraw[black] (B) circle (1.5pt) node[below right] {$M'(b; \alpha)$};
                
                % Nhận xét trực quan dưới hình
                \node[black, fill=white, inner sep=2pt] at (2.5, -1.2) {Phép toán Logarit là ngược của Lũy thừa};
            \end{tikzpicture}
        \end{column}
    \end{columns}
\end{frame}


\subsection{Các quy tắc tính Logarit}

\begin{frame}[fragile]
    \frametitle{\thesection. \insertsection}
    \framesubtitle{\thesection.\thesubsection. Các quy tắc tính (Rules of Logarithms)}

    % Khối lý thuyết điều kiện chung ở trên cùng
    \begin{block}{Điều kiện áp dụng}
        \footnotesize
        Cho số thực dương $a \neq 1$ và các số thực dương $M, N > 0$. Ta có các quy tắc tính logarit sau đây:
    \end{block}
    
    \vspace{0.4cm}
    
    % Dàn ngang 3 khối quy tắc ở dưới để học sinh dễ so sánh
    \begin{center}
    \setlength{\tabcolsep}{6pt} % Chỉnh khoảng cách giữa các cột
    \begin{tabular}{c c c}
        
        % --- QUY TẮC 1: TÍCH ---
        \begin{tikzpicture}
            \node[draw=blue, thick, rounded corners, fill=blue!5, inner sep=10pt, text width=4.5cm, align=center] {
                \footnotesize \textbf{\textcolor{blue}{1. Logarit của một tích}} \\[0.2cm]
                Logarit của một tích bằng \\ \textbf{tổng} các logarit: \\[0.2cm]
                $\boxed{\log_a (M \cdot N) = \log_a M + \log_a N}$
            };
        \end{tikzpicture}
        & 
        
        % --- QUY TẮC 2: THƯƠNG ---
        \begin{tikzpicture}
            \node[draw=red, thick, rounded corners, fill=red!5, inner sep=10pt, text width=4.5cm, align=center] {
                \footnotesize \textbf{\textcolor{red}{2. Logarit của một thương}} \\[0.2cm]
                Logarit của một thương bằng \\ \textbf{hiệu} các logarit: \\[0.2cm]
                $\boxed{\log_a \left(\frac{M}{N}\right) = \log_a M - \log_a N}$
            };
        \end{tikzpicture}
        & 
        
        % --- QUY TẮC 3: LŨY THỪA ---
        \begin{tikzpicture}
            \node[draw=green!60!black, thick, rounded corners, fill=green!5, inner sep=10pt, text width=4.5cm, align=center] {
                \footnotesize \textbf{\textcolor{green!60!black}{3. Logarit của lũy thừa}} \\[0.2cm]
                Đưa số mũ $\alpha$ ra \textbf{ngoài} \\ dấu logarit: \\[0.2cm]
                $\boxed{\log_a (M^\alpha) = \alpha \log_a M}$
            };
        \end{tikzpicture}
        \\
        
        % --- GHI CHÚ BÊN DƯỚI TỪNG KHỐI ---
        \begin{minipage}{4.5cm}
            \vspace{0.2cm}
            \centering \scriptsize \textcolor{gray}{Mở rộng: $\log_a(x_1 x_2...x_n) = \sum \log_a x_i$}
        \end{minipage}
        &
        \begin{minipage}{4.5cm}
            \vspace{0.2cm}
            \centering \scriptsize \textcolor{gray}{Đặc biệt: $\log_a \left(\frac{1}{N}\right) = - \log_a N$}
        \end{minipage}
        &
        \begin{minipage}{4.5cm}
            \vspace{0.2cm}
            \centering \scriptsize \textcolor{gray}{Đặc biệt: $\log_a (\sqrt[n]{M}) = \frac{1}{n} \log_a M$}
        \end{minipage}
        
    \end{tabular}
    \end{center}

\end{frame}

\end{document}